LLM harness 是包裹冻结模型的一组提示、工具与控制流程程序。自动生成大量这类程序带来一个评估问题：什么证据能够说明新程序提供了新能力？不同代码可能执行相同路径，不同路径可能给出相同答案，而一次运行中不同的答案可能来自执行随机性，并非可重复的优势。

我们在 BIRD 文本转 SQL 任务上研究这一问题 \citep{bird2021}。在 Phase-I 发现种群中，66 对生成候选中有 21 对的正确性向量相同；加入预置的 \react{} 和 \bare{} 基线后，为 91 对中的 36 对。代码相似度与结果分歧没有观测到秩关联（$\rho=-0.01$）。轨迹检查区分出三种情况：宣称的机制未进入控制流程；机制存在但在所检查题目上未触发；多轮实现失效。这些观察促使我们审核测量流程，不能证明生成的 harness 普遍坍缩。

后续预先指定的析因研究在六个构建器、三个种子上交叉设置策略强制与一致性门控。主评估使用协议开发之外的九个数据库、1,169 道题，并披露保留了一个 18 题冒烟测试例外。在这一调整后的设置中，无门控自由生成已产生 8.6 个百分点的观测 oracle 提升空间。Phase-I 模式没有再次出现，但生成方式、评分器、数据和执行均有变化，缺少受控桥接实验时不能归因于某一个协议要素。

原门控还改变了允许进入种群的策略空间：两种允许的单次调用提示变换无法通过机制检查。因此，该干预把机制筛选与一致性验证混合在一起。我们保留原实验，报告版本化的审后重分析，不将其视为修复后门控的试验。主对比仍为 $\RevisionPrimaryDelta$ 个百分点（95\% CI $\RevisionPrimaryCI$）；该估计既未证明改善，也未确立等效性。

当前贡献是对三项测量边界的实证说明：
\begin{enumerate}\tightlist
\item \textbf{代码、轨迹和结果多样性不同。}发现阶段的例子说明，未激活的分支与缺失的机制需要不同的检验。
\item \textbf{准入规则和推断方式决定比较含义。}析因研究及其数据、统计链条的明确修正，说明实际评估了哪项干预，以及估计支持什么。
\item \textbf{观测 oracle 价值不等于稳定或可用价值。}重复波动和探索性选择结果要求分别报告 oracle 准确率、最佳固定成员准确率、跨重复优势与实际选择器表现。
\end{enumerate}
决定性问题仍是：真实 harness 种群能否超过相同预算下同一代码的独立执行种群，而且优势能否在独立重复中保留？当前研究尚未完成这一对照。因此，稳定互补性是待检验的估计目标，尚非已经发现的性质。
